\subsection{Evaluation Metrics}
\label{sec:metrics}

Let $y_i$ be a reference GHI value, $\widehat{y}_i$ its post-processed
forecast, $e_i=y_i-\widehat{y}_i$, $\bar{y}$ the mean reference value in the
same evaluation vector, and $n$ its number of target points. Each metric is
computed first for one site, one model estimate, one partition, one declared
horizon, and one horizon scope.

\subsubsection{Mean absolute error}

\begin{equation}
 \mathrm{MAE}
 =
 \frac{1}{n}\sum_{i=1}^{n}\left|e_i\right|.
 \label{eq:mae}
\end{equation}
MAE is the primary error measure and remains in W/m$^2$. Its linear penalty
makes the magnitude directly interpretable and less dominated by individual
large deviations than a squared-error score \cite{hyndman2006}. Lower values
are better.

\subsubsection{Root mean squared error}

\begin{equation}
 \mathrm{RMSE}
 =
 \sqrt{\frac{1}{n}\sum_{i=1}^{n}e_i^2}.
 \label{eq:rmse}
\end{equation}
RMSE is also expressed in W/m$^2$, but squaring gives larger deviations more
weight. Reporting it alongside MAE reveals whether occasional large errors
change the model ordering, a standard concern in deterministic solar-forecast
verification \cite{yang2020verification}. Lower values are better.

\subsubsection{Normalized root mean squared error}

\begin{equation}
 \mathrm{nRMSE}
 =
 \frac{\mathrm{RMSE}}{\bar{y}},
 \qquad
 \mathrm{nRMSE}_{\%}=100\,\mathrm{nRMSE}.
 \label{eq:nrmse}
\end{equation}
The machine-readable artifacts store the dimensionless ratio; tables that use
percent explicitly multiply it by 100. Normalization supports comparison among
sites with different mean irradiance levels, but it inherits the instability
of a small denominator \cite{voyant2022}. The value is therefore left undefined
when $\bar{y}$ is numerically zero. Lower values are better.

\subsubsection{Coefficient of determination}

\begin{equation}
 R^2
 =
 1-
 \frac{\sum_{i=1}^{n}e_i^2}
      {\sum_{i=1}^{n}(y_i-\bar{y})^2}.
 \label{eq:r2}
\end{equation}
$R^2$ measures variation tracked relative to a constant predictor built from
the evaluation mean \cite{yang2020verification}. It is undefined for fewer
than two targets or a constant reference vector and can be negative when the
forecast is worse than that constant reference. Higher values are better;
$R^2$ does not replace the scale-dependent errors.

\subsubsection{Horizon scopes and spatial aggregation}

For a declared horizon $H$, the cumulative scope concatenates leads
$1,\ldots,H$ over all valid origins, whereas the exact-lead scope retains only
lead $H$. The two scopes are identical at $H=1$. This distinction is applied at
24, 48, and 72 hours; 1, 3, 7, 14, and 30 days; and, for the exploratory
multi-month task, 1, 3, and 6 months. A horizon label therefore always denotes
both a numerical lead and an explicitly named scope.

If $M_{\ell}$ is any valid site-level metric and $L=10$, the macro result is
\begin{equation}
 M_{\mathrm{macro}}
 =
 \frac{1}{L}\sum_{\ell=1}^{L}M_{\ell}.
 \label{eq:macro_metric}
\end{equation}
Equal site weighting prevents a site with more valid target points from
dominating the geographical summary. Macro-MAE is the predeclared primary
metric at every resolution; MAE, RMSE, nRMSE, and $R^2$ are nevertheless
retained together rather than selectively reported.

For daily and monthly learned models, seed-level metrics quantify variability,
whereas the TimesNet--DilatedRNN paired comparison uses the arithmetic
five-seed ensemble forecast. The hourly extension compares the single
registered seed 42 for both models. For each site, metric, horizon, and scope,
the pipeline stores
\begin{equation}
 \Delta_{\ell,M}
 =
 M_{\ell,\mathrm{TimesNet}}-M_{\ell,\mathrm{DilatedRNN}}.
 \label{eq:paired_difference}
\end{equation}
A negative $\Delta_{\ell,M}$ is a TimesNet win for the three error metrics; a
positive value is a win for $R^2$. Mean and median differences and the number
of site wins are descriptive. Overlapping multi-step windows are not treated
as independent observations for an i.i.d. significance test.
